\documentclass[12pt]{article}
\setcounter{secnumdepth}{0}

\usepackage[margin=1in]{geometry}
\usepackage{amsmath,amssymb,amsthm}
\usepackage{hyperref}

\setlength{\parindent}{0pt}
\sloppy
\raggedright

\newtheorem*{fact*}{Fact}
\newtheorem*{nomn*}{Nomenclature}
\newtheorem*{fact_plus_nomn*}{Fact + Nomenclature}
\newtheorem*{exer*}{Exercise}
\newtheorem{claim}{Claim}
\newtheorem*{theorem*}{Theorem}
\newtheorem{theorem}{Theorem}
\newtheorem{conjecture}{Conjecture}
\newtheorem{proposition}{Proposition}
\newtheorem{lemma}{Lemma}
\newtheorem{corollary}{Corollary}
\newtheorem{remark}{Remark}
\newtheorem{defn}{Definition}
\newtheorem{example}{Example}

\newcommand{\abs}[1]{\left\lvert#1\right\rvert}
\DeclareMathOperator{\aut}{Aut}
\DeclareMathOperator{\ber}{Ber}
\newcommand{\eps}{\varepsilon}
\DeclareMathOperator{\id}{id}
\renewcommand{\d}{\mathrm{d}}
\newcommand{\gen}[1]{\langle#1\rangle}
\newcommand{\m}{\mathfrak{m}}
\newcommand{\norm}[1]{\left\lVert#1\right\rVert}
\newcommand{\q}{\mathfrak{q}}
\DeclareMathOperator{\spn}{span}
\DeclareMathOperator{\trace}{Tr}

\newcommand{\C}{\mathbf{C}}
\DeclareMathOperator{\End}{End}
\newcommand{\F}{\mathbf{F}}
\newcommand{\Q}{\mathbf{Q}}
\newcommand{\R}{\mathbf{R}}
\newcommand{\Z}{\mathbf{Z}}
\newcommand{\OO}{\mathcal{O}}




\begin{document}


This document is written with the intention of helping the author learn. Not only by teaching, but also in the scenario where a time traveler gifts it to their younger self.
Other readers are also welcome, but be warned: this text is written for highly motivated, mathematically mature readers. Readers are expected to be able to supply rigorous proofs where they are omitted, try their hand at proving results before reading the given proofs, internalize the given proofs in their own terms, work out cases on their own, and so on. 
The author believes the best way to learn is with a mentor. The author knows there are easily available books far better than this document, in every respect, except that it lets the author do as they please.
This document, eventually, should be in guided exercise form. Exposition should follow the historical progress and have natural motivation.

\tableofcontents
 

\newpage
\section{The Fundamental Theorem of Arithmetic}


\begin{fact*}[Principle of Induction]
Any non-empty subset of the positif integers contains a smallest element.
\end{fact*}


\begin{fact_plus_nomn*}
Let $d$ be a positive integer and $m$ be an integer. Then, there exists a unique pair of integers $q,r$ satisfying $m=dq+r$ and $0\le r < d$. We call $q$ the \underline{quotient} and $r$ the \underline{remainder} of the division of $m$ by $d$. When $r=0$, that is, when there exists an integer $q$ with $m=dq$, we say that $d$ \underline{divides} $m$, that $d$ is a \underline{divisor} of $m$, or that $m$ is \underline{divisible} by $d$, and we write this relation as $d\mid m$.

\end{fact_plus_nomn*}

\begin{nomn*}
An integer $p \ge 2$ is called a \underline{prime number} if its only divisors are $1$ and $p$. 
\end{nomn*}


\begin{fact*}
Each positive integer $n\ge 2$ is divisible by at least one prime number $p$.
\end{fact*}


\begin{fact*}[Primes don't materialize out of thin air]
Let $p$ be a prime number, and let $m,n$ be integers satisfying $p\mid mn$. Then, either $p\mid m$ or $p\mid n$.
\end{fact*}


\begin{proof}
Without loss of generality, we may restict our attention to positif $m,n$.
By the principle of induction, we may suppose the theorem holds for all prime numbers $p'$ smaller than $p$ (and any $n,m$). Again by the principle of induction, we may suppose the theorem holds for all positif pairs $m',n'$ with a smaller sum $m' + n' < m + n$ (and our fixed $p$).
Applying division with remainder, we may suppose $0 \le m,n < p$. 
We may now write $kp=mn$ for some integer $0 \le k < p$. \\
Case 1. If $k = 0$ then either $m=0$ or $n=0$, and we have our desired conclusion. \\
Case 2. We cannot have $k=1$, by $p$ being a prime number. \\
Case 3. If $k\ge 2$, then we may pick some prime number $q$ dividing $k$. In particular $q \mid mn$ and $q < p$. Thus, the (first) induction hypothesis, yields that either $q\mid m$ or $q\mid n$. Without loss of generality, we may assume $q\mid m$. We may write $k=qd$ and $m=qr$ for some positif integers $d,r$ with $r < m$, yielding $dp=rn$. Thus, the (second) induction hypothesis yields that either $p\mid r$ (implying $p\mid m$) or $p\mid n$, reaching our desired conclusion.
\end{proof}


\begin{fact*}[Fundamental Theorem of Arithmetic]
Each positive integer $n$ may be written as a product of a list of prime numbers $n=p_1\dots p_\ell$, in a manner unique up to reordering of the terms.
\end{fact*}


\newpage
\section{Number Theory - Exercises for Lecture 1}


\begin{fact*}
$d \mid m$ and $m \mid n$ implies $d\mid n$.
\end{fact*}


\begin{fact*}
$m$ and $n$ have the same remainder with respect to division by $d$ if and only if $d \mid m-n$.    
\end{fact*}


\begin{fact*}
The remainders of $m+n$ and of $mn$ with respect to division by $d$ depends only on the remainders of $m$ and of $n$ with respect to division by $d$. 
\end{fact*}


\begin{fact_plus_nomn*}
For a positive integer $n$ and a prime number $p$, let the \underline{$p$-adic valuation} $\nu_p(n)$ be the number of times we may divide $n$ by $p$ (until the result is no longer divisible by $p$). 
We have \[ n = \prod_{p} p^{\nu_p(n)} \]
\end{fact_plus_nomn*}


\begin{fact*}
We have $\nu_p(nm) = \nu_p(n) + \nu_p(m)$.
\end{fact*}


\begin{fact*}
We have $\nu_p(n + m) \ge \min(\nu_p(n), \nu_p(m))$, and equality holds if $\nu_p(n) \neq \nu_p(m)$ (but not necessarily otherwise).
\end{fact*}


\begin{fact*}
A positive integer is a perfect $k$-th power if and only if $k\mid \nu_p(n)$ for all prime numbers $p$. Deduce that the square root of $2$ is irrational. More generally, the $k$-th root of a positive integer which is not a perfect $k$-th power, is irrational.
\end{fact*}


\begin{fact*}
The number of divisors of $n$ equals $\prod_p (\nu_p(n) + 1)$. In particular, $n$ has an odd number of divisors if and only if $n$ is a perfect square - convince a child of this.
\end{fact*}


\begin{exer*}
Search for Zarmelo's proof of the Fundamental Theorem of Arithmetic, which does not make use of the fact $p\mid mn\implies p\mid m$ or $p \mid n$.
\end{exer*}


\newpage


TODO: fix lecture 1 given the following.\\


\begin{nomn*}
We call $a$ \underline{invertible} modulo $n$ if $0,a,2a,3a,\dots,(n-1)a$ cover all $n$ distinct residue classes modulo $n$. Namely, if the row corresponding to $a$ in the multiplication table modulo $n$ contains each residue precisely once. 
\end{nomn*}


\begin{fact*}
$a$ is invertible modulo $n$ if and only if there exists $b$ such that $ab\equiv 1\mod n$. In that case, $b$ is unique modulo $n$.
\end{fact*}


\begin{fact*}
$xy$ is invertible modulo $n$ if and only if both $x$ and $y$ are invertible modulo $n$.
\end{fact*}


\begin{fact*}
Let $p$ be a prime number. Then each non-zero residue $a$ is invertible modulo $p$.
\end{fact*}


\begin{proof}
By induction on $a=1,\dots,p-1$. For $a=1$ the result is clear. Assume the result holds for all previous $a$ values. Divide $p$ by $a$ with remainder so that $p=aq+r$ for some $0\le r < a$. Since $p$ is a prime number and $1 < a < p$, we have that $a\nmid p$ so $1\le r < a$, and thus the inductive hypothesis yields that $r\equiv -aq \mod p$ is invertible modulo $p$, and hence so is $a$. 
\end{proof}


\begin{nomn*}
The number of invertible residue classes modulo $n$ is denoted $\varphi(n)$.
\end{nomn*}


\begin{fact*}
The multiplication table modulo $n$, reduced to the rows and columns of invertible residues modulo $n$ (so that it has size $\varphi(n)\times \varphi(n)$), has only entries that are invertible modulo $n$, and moreover each row (and column) contains each invertible residues precisely once. 
\end{fact*}


\begin{fact*}
For each invertible residue class $a$ modulo $n$ we have $a^{\varphi(n)}\equiv 1\mod n$.
\end{fact*}


\begin{proof}
Let $x_1,\dots,x_{\varphi(n)}$ list the invertible residue classes modulo $n$. Then $ax_1,\dots,ax_{\varphi(n)}$ is a reordering of the same list. Therefore, the products of the two lists are equal, that is $x_1\dots x_{\varphi(n)} \equiv ax_1 \dots ax_{\varphi(n)} = a^{\varphi(n)}x_1\dots x_{\varphi(n)} \mod n$. Since $x_1,\dots,x_{\varphi(n)}$ are invertible modulo $n$, we have $1\equiv a^{\varphi(n)} \mod n$ as desired.
\end{proof}


\newpage
\section{Quadratic Residues and Quadratic Forms}

% taken from book Modern Olympiad Number Theory
\begin{fact*}[Thue]
Let $a$ be coprime to $n$, and $n\ge 2$. Then, there exists a solution to $ax\equiv y\mod n$ with $0 < \abs{x}, \abs{y} < \sqrt{n}$.
\end{fact*}


\begin{proof}
Consider all combinations $ax-y\mod n$ where $x,y=0,\dots,\lfloor \sqrt{n} \rfloor$. There are $(1+\lfloor \sqrt{n} \rfloor)^2>n$ such combinations, so by the pigeon hole principle there must be a repetition. So let $ax_1-y_1\equiv ax_2-y_2$ where $(x_1,y_1)\neq(x_2,y_2)$. Since $a$ is coprime to $n$, we must have that $x_1\not\equiv x_2$ and $y_1\not\equiv y_2$. Now $x=x_1-x_2$ and $y=y_1-y_2$ satisfy the requested requirments, as desired.
\end{proof}


\begin{fact*}
Let $p\equiv 1\mod 4$. Then, there exists $x,y$ with $p=x^2+y^2$.
\end{fact*}


\begin{proof}
Let $a$ solve $a^2\equiv -1\mod p$. Applying Thue's result, let $x,y$ solve $ax\equiv y\mod p$ with $0 < \abs{x},\abs{y} < \sqrt{p}$. Then $0 < x^2+y^2 < 2p$ and $x^2+y^2\equiv x^2(1+a^2)\equiv 0\mod p$, so $p=x^2+y^2$ as desired.
\end{proof}


\newpage
\section{Algebraic Integers}


\begin{fact*}
A number which is both rational and an algebraic integer must be an integer.
\[\Q \cap \overline\Z = \Z\]
\end{fact*}


\begin{fact*}[Fundamental Theorem of Symmetric Polynomials]
Let $f(x_1,\dots,x_n)$ be a symmetric polynomial. Then, in fact, $f(x_1,\dots,x_n) = g(e_1,\dots,e_n)$ may be given as a polynomial expression in the elementary symmetric polynomials in $x_1,\dots,x_n$. 
\end{fact*}


\begin{proof}
We give a lexicographic ordering of monomials, so that $x_1^{i_1}\dots x_n^{i_n} >_{lex} x_1^{j_1}\dots x_n^{j_n}$ whenever $i_k > j_k$ for the first index $k$ where $i_k\neq j_k$. This is a total order of all monomials. Suppose the lexicographically highest monomial term of $f(x_1,\dots,x_n)$ is $x_I=x_1^{i_1}\dots x_n^{i_n}$, and that the theorem holds for all symmetric polynomials containing terms only lexicographically smaller than this monomial. Since $f$ is symmetric, we have $i_1\ge \dots \ge i_n$. Now, $E_I=e_1^{i_1-i_2}e_2^{i_2-i_3}\dots e_{n-1}^{i_{n-1}-i_n}e_n^{i_n}$ has highest lexicographic monomial term equal to $x_1^{i_1-i_2}(x_1x_2)^{i_2-i_3}\dots(x_1\dots x_{n-1})^{i_{n-1}-i_n}(x_1\dots x_n)^{i_n} = x_I$. Subtracting $cE_I$ from $f$ (where $c$ is the coefficient of $x_I$ in $f$) yields a symmetric polynomial, all of whose terms are lexicographically smaller than $x_I$, so by induction $f-cE_I$ is expressible as a polynomial in $e_1,\dots,e_n$, and thus so is $f$, as desired.
\end{proof}


\begin{fact*}
Sums and products of algebraic integers are algebraic integers.
\end{fact*}


\begin{proof}
Let $\alpha,\beta$ be algebraic integers. Then we may write
\[
(x-\alpha_1)\dots (x-\alpha_m) = x^m + a_1x^{m-1} + \dots + a_{m-1}x + a_m
\]
\[
(x-\beta_1)\dots (x-\beta_n) = x^n + b_1x^{n-1} + \dots + b_{n-1}x + b_n
\]
for integers $a_i,b_j$ and numbers $\alpha_i,\beta_j$ with $\alpha=\alpha_1$, $\beta=\beta_1$. Consider the polynomials
\[
\prod_{i,j}(x-(\alpha_i+\beta_j))
\]
and
\[
\prod_{i,j}(x-(\alpha_i\beta_j))
\]
By the fundamental theorem of symmetric polynomials, each of these is expressible as a polynomial in $x$ whose coefficients are polynomials in $a_i,b_j$. That is, these polynomials have integer coefficients. In particular, $\alpha+\beta$ and $\alpha\beta$ are algebraic integers, as desired.
\end{proof}


\begin{proof}
Let $A\in\Z^{m\times m}$, $B\in\Z^{n\times n}$ be respectively the companion matrices of $\alpha$, $\beta$.
Let $A$ have eigenvalues $\alpha_i$ with corresponding eigenvectors $v_i$, and let $B$ have eigenvalues $\beta_j$ with corresponding eigenvectors $u_j$.
Then $A\otimes B\in \Z^{nm\times nm}$ and $A\otimes \id_n + \id_m\otimes B$ have respectively the eigenvalues $\alpha_i\beta_j$ and $\alpha_i+\beta_j$, both with corresponding eigenvectors $v_i\otimes u_j$. It follows that $\alpha_i\beta_j$ and $\alpha_i+\beta_j$ (and in particular $\alpha\beta$ and $\alpha+\beta$) are algebraic integers, as desired.
\end{proof}


\newpage
\section{Cyclotomics}


\begin{nomn*}
The $n$-th cyclotomic polynomial is
\[
\Phi_n(x) = \prod_{\mu\text{ a primitive $n$-th root of unity}} (x-\mu) = \prod_{j\in(\Z/n\Z)^\times} (x-\zeta^{j})
\]
where $\zeta=e^{2\pi i/n}$.
\end{nomn*}


\begin{fact*}
\[
x^n - 1 = \prod_{d\mid n} \Phi_d(x)
\]
\end{fact*}


\begin{fact*}
$\Phi_n(x)\in\Z[x]$ is a monic integer polynomial.
\end{fact*}


\begin{proof}
Clearly $\Phi_n(x)$ is a monic polynomial. By induction, we have that $\Phi_n(x)$ is a rational function of $x$ with rational coefficients, but it has no poles so it is a rational polynomial. $\Phi_n(x)\in\Q(x)\cap \C[x] = \Q[x]$. It is clear that the roots of $\Phi_n(x)$ are algebraic integers, so the coefficients of $\Phi_n(x)$ are algebraic integers that are rational numbers, and so $\Phi_n(x)$ has integer coefficients.
$\Phi_n(x)\in\Q[x]\cap \overline\Z[x] = \Z[x]$.
\end{proof}


\begin{fact*}
$\Phi_n(x)$ is irreducible in $\Q[x]$.
\end{fact*}


\begin{proof}(Schur)
Recall that the discriminant of $x^n-1$ is $\pm n^n$. Let $f(x)$ be the minimal polynomial of $\zeta$ over $\Q$. So $f(x)\in\Z[x]$ and $f(x)=(x-\zeta_1)\dots(x-\zeta_k)$ for some distinct $n$-th roots of unity $\zeta_1,\dots,\zeta_k$. We must show that the roots of $f(x)$ are closed under exponentiating by powers coprime to $n$. It suffices to show that if $\mu$ is a root of $f$ and $p$ is a prime coprime to $n$ then $\mu^p$ is a root of $f$. We have that $f(\mu^p) = f(\mu^p) - f(\mu)^p = f(x^p) - f(x)^p \restriction_{x=\mu} = pg(x) \restriction_{x=\mu}$ for some $g(x)\in\Z[x]$, so that $p\mid f(\mu^p)$ over $\overline\Z$. Moreover, if we assume for the sake of contradiction that $\mu^p$ is not a root of $f$ then $f(\mu^p) = (\mu^p-\zeta_1)\dots(\mu^p-\zeta_k)$ divides the discriminant $\pm n^n$ over $\overline\Z$. However, $p\mid \pm n^n$ over $\overline\Z$ implies $p\mid n$ over $\Z$, contrary to our assumption. This shows that $\mu^p$ is a root of $f$, as desired.
\end{proof}


\begin{fact*}
Let $\zeta=\zeta_{p^r}$, $K=\Q(\zeta)$, and $m=[K:\Q]$. We have 
\begin{itemize}
    \item $N_{K/\Q}(1-\zeta)=p$
    
    \item $\gen{1-\zeta}$ is a prime ideal, lying over $p$. $(1-\zeta)\OO_K \cap \Z = p\Z$

    \item $(1-\zeta^j)/(1-\zeta)\in\OO_K^\times$ for all $j\in(\Z/p^r\Z)^{\times}$
    
    \item $p\OO_K = (1-\zeta)^m \OO_K$
    
    \item $\Z[\zeta] \cap p\OO_K = p\Z[\zeta]$
    
    \item $\OO_K = \Z[\zeta]$
\end{itemize}
\end{fact*} % this is taken from https://virtualmath1.stanford.edu/~conrad/154Page/handouts/cycint.pdf where they use it to show O_K=Z[zeta] using discriminants
% but perhaps we can also use traces for this?


\newpage
\section{Fermat's Last Theorem}
\begin{theorem*}(Wiles, Taylor, Taniyama, Shimura, Ribet, Frey, \dots, Fermat)
Let $n\ge 3$. Then, any solution in rational numbers to the equation \[x^n+y^n=z^n\]
is trivial, in the sense that $xyz=0$.
\end{theorem*}


\begin{exer*}
It suffices to prove the case $n=4$ and the cases where $n=p$ is an odd prime. Also, we may assume $x,y,z$ are pairwise relatively prime.
\end{exer*}


\begin{exer*}
There are algebraic integers on the unit circle other than the roots of unity.
\end{exer*}


% \begin{proof}
% Let $f(z)=z^4-2z^3-2z+1$. $f$ has two real roots and a pair of conjugate roots (calculus). But $f$ has symmetric coefficients, meaning its roots are closed under taking inverse. So $f$ has two roots on the unit circle. $f$ is irreducible over $\Q$ (for example, $f(x+1)$ is Eisenstein), and is distinct from any cyclotomic polynomial (recall that $\phi(n)\ge \sqrt{n}$ for $n>6$).
% \end{proof}


However, such numbers have conjugates outside the unit circle.


\begin{fact*}
Suppose $f(x)=\prod_\alpha(x-\alpha)\in\Z[x]$ is a monic polynomial with integer coefficients, all of whose roots lie in the closed unit circle. Then each roots $\alpha$ is either zero or a root of unity.
\end{fact*}


\begin{proof}
Without loss of generality $f(x)$ is irreducible, so it has no repeated roots. Let $A$ be the companion matrix. Then $A$ is diagonable, with eigenvalues all of absolute value at most one. Thus the matrices $(A^n)_{n\ge 1}$ are bounded. However, these are integer matrices. By the pigeonhole principle, we have $A^n = A^{n+k}$ for some $k>0$. It follows that $\alpha^n=\alpha^{n+k}$, so either $\alpha=0$ or $\alpha$ is a $k$-th root of unity.
\end{proof}


\begin{fact*}
Any unit $u$ of $\Q(\zeta_p)$ may be written as $u=r\zeta^k$ for some real $r$ and some integer $k$.
\end{fact*}


\begin{fact*}
For any $\alpha\in\Z[\zeta_p]$ there is an integer $a$ with $\alpha^p \equiv a \mod \gen{1-\zeta}^p$. 
\end{fact*}


\begin{fact*}[Kummer]
Let $p$ be a regular prime. Then the $n=p$ case of FLT is true.
\end{fact*}


\begin{proof}
Let us write the equation in the form \[x^p+y^p+z^p=0\]
We have 
\[\prod_{j=0}^{p-1}(x+\zeta^j y) = x^p + y^p = -z^p\]
And so 
\[\prod_{j=0}^{p-1}\gen{x+\zeta^j y} = \gen{z}^p\]
First let us handle the case $p\nmid xyz$. \\
Claim. We claim that $\gen{x+\zeta^j y}$ are pairwise coprime. \\
Indeed, suppose $\q \mid \gen{x+\zeta^j y}$, $\q \mid \gen{x+\zeta^k y}$. Then $x+\zeta^j y, x+\zeta^k y\in \q$. Thus $ \zeta^k(1-\zeta^{k-j})y\in\q$ $\implies$ $y\in\q$ or $1-\zeta^t \in \q$. However, if $y\in\q$ then $x\in\q$, but this contradicts $x,y$ being coprime. It follows that $1-\zeta^t \in\q$. Now $1-\zeta^t \sim 1-\zeta$, and $\gen{1-\zeta}$ is prime, so $\q = \gen{1-\zeta}$. Thus $z\in\gen{1-\zeta}$. Since $N(\gen{1-\zeta})=p$, we get $p\mid z$, contradicting our hypothesis. \\
It follows that each of the ideals $\gen{x+\zeta^j y}$ is a $p$-th power. Now let $\gen{x+\zeta y} = \m^p$. So $\m$ has order dividing $p$ in the class group. Since $p$ is regular, $\m=\gen{\alpha}$ must be principal. We have \[x+\zeta y = u\alpha^p\]
for some unit $u$. We may write $u=r\zeta^k$ for some real $r$ and integer $k$. We find an integer $a$ with $\alpha^p\equiv a\mod \gen{1-\zeta}^p$ so $\alpha^p \equiv a\mod p$ and hence $x+\zeta y \equiv r\zeta^k a \mod p$. Inverting by $\zeta^k$ and conjugating, we find
\[\zeta^{-k}(x+\zeta y) - \zeta^{k}(x+\zeta^{-1} y) \equiv 0 \mod p\]
So consider the algebraic integer $\beta$ with 
$\beta = p^{-1}\left(\zeta^{-k}x + \zeta^{1-k}y - \zeta^{k}x - \zeta^{k-1}y\right)$. Since $\OO_K = \Z[\zeta]$ has integer basis ($1,\zeta,\dots,\zeta^{p-1}$ minus any particular one), if $\zeta^{-k},\zeta^{1-k},\zeta^{k},\zeta^{k-1}$ were all distinct we'd get that $p\mid x,y$, contradicting our hypothesis. \\
Case 1. We have $k\equiv 0\mod p$. In this case $\gen{1-\zeta}^{p-1}=\gen{p} \ni (\zeta-\zeta^-1)y = \zeta^{-1}(\zeta^2-1)y$. Since $\zeta^2-1\sim 1-\zeta$, we have $y\in\gen{1-\zeta} \cap \Z = p\Z$, contradicting our hypothesis. \\
Case 2. We have $k\equiv 1\mod p$. This is the same as the previous case with $y$ replaced by $x$. \\
Case 3. We must have $2k\equiv 1\mod p$. Thus \[p\beta \zeta^k = (x-y)(1-\zeta)\]
Taking norms, we have $p\mid x-y$. By symmetry, $x\equiv y\equiv z \mod p$ and so $0\equiv 3x^p \mod p$. Our assumption yields $p=3$. However, for $3\nmid x,y,z$ we have $0 = x^3+y^3+z^3 \equiv \pm1\pm1\pm1 \mod 9$, a contradiction. \\

TODO handle case $p\mid xyz$ and prove lemmas. % https://wstein.org/129-05/final_papers/Emily_Riehl.pdf
% https://kconrad.math.uconn.edu/blurbs/gradnumthy/kummer.pdf
\end{proof}


\newpage
\section{Gauss Sums and Quadratic Reciprocity}


Let $p$ be an odd prime, $\eps_p=(-1)^{\frac{p-1}{2}}$, $p^*=\eps_p p$, $\zeta = e^{2\pi i/p}$. Now, we set
\[\tau = \sum_{x\in\F_p}\zeta^{x^2}\]
We have
\[\tau^* = \sum_{y\in\F_p}\zeta^{-y^2}\]
If $\eps_p=1$ then $-1$ is a square modulo $p$ and so $\tau^*=\tau$. Otherwise $-1$ is a non-square modulo $p$ and so $\tau+\tau^*=2\sum_{w\in\F_p}\zeta^{w}=0$. 
Thus \[\tau^*=\eps_p\tau\]
We have
\[\tau\tau^* = \sum_{x,y\in\F_p}\zeta^{x^2-y^2}=\sum_{u,v\in\F_p}\zeta^{uv} = p\]
where we make use of the bijection $(x,y)\longleftrightarrow(u,v)$ given by $u=x+y$, $v=x-y$.
It follows that
\[\tau = \pm \sqrt{p^*}\]
%The sign is $+$, but this is not obvious. TODO prove.


\begin{fact*}
We have
\[\zeta = \sum_{a\in\F_p}(a:p)\zeta^a\]
\end{fact*}


\begin{fact*}
Let $a_0,\dots,a_{p-1}$ be integers with $\sum_{j}a_j\zeta^j = 0$. Then $a_0=\dots=a_{p-1}$ are all equal.
\end{fact*}


\begin{defn}
Let $N_r(a) = \#\{(x_1,\dots,x_r)\in\F_p^r : x_1^2+\dots+x_r^2=a\}$.
\end{defn}


For example, we have $N_1(a)=1+(a:p)$.


Let $r$ be even. Then 
\[(p^*)^{r/2}=\tau^r = \sum_{x_1,\dots,x_r\in\F_p}\zeta^{x_1^2+\dots+x_r^2} = \sum_{a\in\F_p}N_r(a)\zeta^{a}\]
So
\[N_r(0)-(p^*)^{r/2} = N_r(a) \phantom{-} \forall a\in\F_p^*\]
It follows (since $\sum_a N_r(a)=p^r$) that
\[N_r(a) = \begin{cases}
    p^{-1}(p^r-(p^*)^{r/2}) & a\neq 0 \\
    p^{-1}(p^r+(p-1)(p^*)^{r/2}) & a = 0 \\
\end{cases}\]
Let $r$ be odd. Then
\[(p^*)^{(r-1)/2}\sum_{a\in\F_p}(a:p)\zeta^{a}=(p^*)^{(r-1)/2}\tau=\tau^r = \sum_{x_1,\dots,x_r\in\F_p}\zeta^{x_1^2+\dots+x_r^2} = \sum_{a\in\F_p}N_r(a)\zeta^{a}\]
So that $N_r(a) - (p^*)^{(r-1)/2}(a:p) \equiv c$ is a constant. It follows that $c=p^{r-1}$ and so
\[
N_r(a) = p^{r-1} + (p^*)^{(r-1)/2}(a:p)
\]


\begin{fact*}[Quadratic Reciprocity]
Let $p,q$ be distinct odd primes. Then $(p:q)(q:p)=\eps_p\eps_q$.
\end{fact*}


\begin{proof}
Let us consider $N_q(1) \mod q$. On one hand, $N_q(1)=p^{q-1} + (p^*)^{(q-1)/2} \equiv 1 + (p^*:q) \mod q$. On the other hand, the sphere $\{(x_1,\dots,x_q)\in\F_p^q : x_1^2+\dots+x_q^2=1\}$ is acted on by $C_q$ via cyclic rotations. It follows that $N_q(1) \equiv 1 + (q:p) \mod q$.
\end{proof}


\newpage
\section{Beta Function}
\begin{fact*}
\[\int_0^1 p^k (1-p)^m = \dfrac{k!m!}{(k+m+1)!}\]
\end{fact*}


\begin{proof}
Consider the following stochastic process. First, generate $p\sim U[0,1]$. Then, generate $n$ independent Bernoulli$(p)$ trials $b_1,\dots,b_n \sim \ber(p)$. What is the probability of having precisely $k$ successes? The answer, of course, is
\[
\int_0^1 \binom{n}{k}p^{k}(1-p)^{n-k} \d p
\]
However, instead of generating the $b_i$ as Bernoulli trials out of thin air, let us generate $x_1,\dots,x_n\in U[0,1]$ and set $b_i = [x_i < p]$. Then, the questions asks what is the probability that precisely $k$ of the $n$ numbers $x_1,\dots,x_n$ come before $p$. Clearly, the ordering of $p,x_1,\dots,x_n$ is uniform over all permutations of $n+1$ terms, and so this probability is $\frac{1}{n+1}$ for each $k=0,\dots,n$.
\end{proof}


\newpage
\section{Finite Fields}


\begin{fact*}
The field extension $\F_{q^m}/\F_q$ is cyclic; The Galois group is generated by the Frobenius automorphism $x\mapsto x^{q}$.
\end{fact*}

\newpage
\section{Representation Theory}


The Classical approach.


\begin{fact*}[Weierstrass]
Let $k$ be an algebraically closed field, and let $A$ be a commutative $k$-algebra of finite dimension $n$ over $k$, with no non-zero nilpotents. Then $A \cong k^n$.
\end{fact*}


\begin{proof}
By induction on $n$. The case $n=1$ is clear. Assume $n > 1$ and let $\zeta \in A \setminus k$. Let $m(x)\in k[x]$ be the minimal polynomial of $\zeta$. Write $m(x)=\prod_j (x-r_j)^{a_j}$ with $r_j\in k$ and $a_j \ge 1$. Then $\prod_j(\zeta -r_j)$ is nilpotent. It follows that $a_j = 1$ for all $j$. Thus \[ k[\zeta] \cong k[x]/\gen{m(x)} \cong \prod_j k[x]/\gen{x-r_j} \cong k^{d} \]
where $d > 1$ is the degree of $m(x)$. It follows that we can find orthogonal idempotents $e_1,\dots,e_d$ spanning $k[\zeta]$. Thus $A$ decomposes as \[A = e_1A \oplus \dots \oplus e_dA\]
and we're done by induction.
\end{proof}

Thus we'd like a criterion ruling out the existence of non-zero nilpotents in our commutative algebra $A$. Now, if $x$ is nilpotent, then so is $xy$, and so is the multiplication $m_{xy}\in\End_k(A)$, and thus $\trace m_{xy} = 0$ for all $y\in A$. We notice that $\gen{x,y} = \trace m_{x,y}$ is a symmetric bilinear form $A\times A\to k$. 


\begin{fact*}
Let $k$ be field, and let $A$ be a commutative $k$-algebra of finite dimension $n$ over $k$. If the form $\gen{x,y} = \trace m_{xy}$ is non-degenerate, then $A$ has no non-zero nilpotents. (And the reverse holds in characteristic zero).
\end{fact*}


\begin{fact*}[Frobenius]
Let $G$ be a finite group with $s$ conjugacy classes. Then $Z(\C[G]) \cong \C^{s}$.
\end{fact*}


\begin{proof}
By Weierstrass' theorem, we need only show $Z(\C[G])$ has no non-zero nilpotents. 
By the trace condition, we need only show that $\gen{x,y} = \trace m_{xy}$ is non-degenerate. 
So we shall show that $\det \left(\gen{e_i, e_j}\right)_{ij} \neq 0$ for a basis of $\Z(\C[G])$. 
Let us write $C_1,\dots,C_s$ for the conjugacy classes of $G$.
For a class $C_\gamma$, let $e_\gamma = \sum_{g\in C_\gamma} g$. 
Then $(e_\gamma)_{\gamma=1}^{s}$ forms a basis of $Z(C[G])$. 
We have 
\[ e_k e_j = e_j e_k = \sum_i a_{ijk} e_i\]
where 
\[a_{ijk} = \#\{(b,c)\in C_j\times C_k : bc = a\}\]
for any fixed $a\in C_i$. Letting $M_\gamma$ be the $s\times s$ matrix representing multiplication by $e_\gamma$, we have 
\[\left(M_\gamma\right)_{\alpha,\beta} = a_{\alpha\beta\gamma}\]
and thus
\[\gen{e_\alpha,e_\beta} = \trace M_{\alpha}M_\beta = \sum_{\gamma,\delta} a_{\gamma\delta \alpha}a_{\delta\gamma\beta}\]
Now, we let
\[ h_{ijk} = \#\{(a,b,c)\in C_i\times C_j\times C_k : abc = e\} \]
Since $abc=e\iff bca=e$, we have $h_{ijk}=h_{jki}$. Since $abc=e\iff b(b^{-1}ab)c=e$, we have $h_{ijk}=h_{jik}$. Thus $h$ is symmetric in its variables.
For a conjugacy class $C_i$, we write $C_{i'}$ for the inverse class $\{c^{-1} : c\in C_i\}$. Since $abc=e\iff c^{-1}b^{-1}a^{-1}=e$, we have $h_{ijk}=h_{i'j'k'}$. Let us write \[h_i = \# C_i\]
so that $h_i=h_{i'}$ and
\[ a_{ijk}h_i = h_{i'jk} \]
We are now in position to verify the trace condition $\det \gen{e_\alpha,e_\beta} \neq 0$. Permuting the columns, we swap $\beta\leftrightarrow \beta'$. 
Let
\[ P_{\alpha \beta} := \gen{e_\alpha, e_{\beta'}} = \sum_{\gamma,\delta} a_{\gamma\delta\alpha}a_{\delta\gamma\beta'} = \sum_{\gamma,\delta} \frac{h_{\gamma'\delta\alpha} h_{\delta'\gamma\beta'}}{h_\gamma h_\delta} = \sum_{\gamma,\delta} \frac{h_{\gamma\delta\alpha} h_{\gamma\delta\beta}}{h_\gamma h_\delta}\]
where in the last step we used $\gamma\leftrightarrow \gamma'$. 
Let $Q$ be the real $s\times s^2$ matrix 
\[ Q_{\alpha,(\gamma,\delta)} = 
\frac{h_{\gamma\delta\alpha}}{{\sqrt{h_\gamma h_\delta}}}\]
so that \[P=QQ^t\]
So showing that $P$ is nonsingular reduces to showing that $Q$ has rank $s$. Indeed, let $\gamma_0$ be such that $C_{\gamma_0} = \{e\}$. Then the submatrix $Q_{\alpha,(\gamma_0,\delta)} = \sqrt{h_\delta} [\alpha=\delta']$
is non-singular.
\end{proof}


This allows us to define the irreducible characters of $G$ via the orthogonal idempotents of $Z(\C[G])$, and (I should check) show properties like orthogonality. Only later did Frobenius come up with the notion of a representation.


\begin{exer*}
Find the idempotents for $G=S_3$, the first non-abelian group.
\end{exer*}

\newpage
\section{Absolute Values and Valuations}


\begin{nomn*}
Let $K$ be a field. An \underline{absolute value} on $K$ is a function $\abs{\cdot} : K\to \R_{\ge0}$ satisfying the properties
\begin{enumerate}
    \item $\abs{x} = 0 \iff x = 0$
    \item $\abs{xy} = \abs{x}\abs{y}$
    \item $\abs{x+y} \le \abs{x} + \abs{y}$
\end{enumerate}
for all $x,y\in K$. \\
If $\abs{x+y} \le \max(\abs{x},\abs{y})$ for all $x,y\in K$, we call $\abs{\cdot}$ a \underline{valuation}. 
\end{nomn*}


We have $\abs{\pm 1} = 1$ (and in fact $\abs{\zeta}=1$ for any root of unity $\zeta$). We have $\abs{x^{-1}} = \abs{x}^{-1}$. An absolute value equips $K$ with the metric $d_{\abs{\cdot}}(x,y) = \abs{x-y}$. The trivial absolute value is $\abs{x} = [x \neq 0]$. Raising an absolute value to a power $\lambda \in (0,1]$ yields an absolute value. Raising a valuation to a power $\lambda\in(0,\infty)$ yields a valuation.


\begin{fact*}
Two absolute values $\abs{\cdot}_1, \abs{\cdot}_2$ enrich $K$ with the same topology if and only if there exists some $\lambda > 0$ with $\abs{x}_2 = \abs{x}_1^\lambda$ for all $x\in K$.
\end{fact*}


\begin{proof}
The if direction is clear, so we show the only if direction. The set $\{x:\abs{x}_1 < 1\}$ is the set $\{x : \lim_n x^n = 0\}$, so it is determined by the topology. Thus the topology also determines the inverse set $\{x : \abs{x}_1 > 1\}$, as well as $\{x:\abs{x} = 1\}$. We may assume $\abs{\cdot}$ is non-trivial, and fix $y$ with $a=\abs{y}_1 > 1$ and $b=\abs{y}_2 > 1$. Let $x\in K^\times$, and write $\abs{x}_1 = a^\alpha$ for some $\alpha\in\R$. If $q=m/n$ is a rational number with $q<\alpha$, then $\abs{x}_1 > a^q \implies \abs{x^n}_1 > a^m = \abs{y^m}_1\implies \abs{x^n y^{-m}}_1 > 1 \implies \abs{x^ny^{-m}}_2 > 1 \implies \abs{x}_2 > b^q$.
Similarly, we have that $\abs{x}_2 < b^q$ for all rational numbers $q$ with $q>\alpha$. Thus $\abs{x}_2 = b^\alpha = \abs{x}_1^{\lambda}$ where $\lambda=\log b / \log a$.
\end{proof}


\begin{fact*}
Let $\abs{\cdot}$ be a valuation. Then we have $\abs{x}\neq\abs{y} \implies \abs{x+y} = \max(\abs{x},\abs{y})$.
\end{fact*}


\begin{fact*}
Let $\abs{\cdot}$ be an absolute value on $K$. Then $\abs{\cdot}$ is a valuation if and only if the integers are bounded (meaning $\abs{n} \le C$ for all integers $n$ and some constant $C$).
\end{fact*}


\begin{proof}
The only if direction is clear, so we show the if direction. Let $x,y\in K$. Then \[\abs{x+y}^n = \abs{(x+y)^n} = \abs{\sum_{k}\binom{n}{k}x^k y^{n-k}} \le \sum_k \abs{\binom{n}{k}}\abs{x}^k\abs{y}^{n-k} \le (n+1)C\max(\abs{x},\abs{y})^n\]
taking $n$-th root and a limit we get $\abs{x+y}\le \max(\abs{x},\abs{y})$.
\end{proof}


\begin{fact*}
Let $\abs{\cdot}$ be a valuation, and let $r>0$. Then any point 
\end{fact*}


\begin{nomn*}
Let $K$ be a field, let $\abs{\cdot}$ be an absolute value on $K$, and let $V$ be a vector space over $K$. A \underline{norm} on $V$ is a function $\norm{\cdot}:V\to\R_{\ge0}$ satisfying the properties
\begin{enumerate}
    \item $\norm{v}=0 \iff v = 0$
    \item $\norm{\alpha v} = \abs{\alpha}\norm{v}$
    \item $\norm{u+v} \le \norm{u} + \norm{v}$
\end{enumerate}
for all vectors $u,v\in V$ and scalars $\alpha\in K$.
\end{nomn*}


A pair of norms $\norm{\cdot}_1,\norm{\cdot}_2$ on $V$ are called \underline{equivalent} if there are constants $c,C>0$ such that \[ c\norm{v}_2 \le \norm{v}_1 \le C\norm{v}_2\]
for all $v\in V$.

 
\begin{fact*}
Assume $K$ is complete and $V$ is finite dimensional. Then, any two norms on $V$ are equivalent, and any norm is complete.
\end{fact*}


\begin{proof}
By induction on $d=\dim V$. The cases $d=0,1$ are clear, so let $d\ge 2$. Let $\norm{\cdot}$ be any norm on $V$. Let us fix a basis $e_1,\dots,e_d$ of $V$, and define
\[\norm{v}_\infty = \max(\abs{x_i})\]
where $v=x_1e_1+\dots+x_de_d$. Then $\norm{\cdot}_\infty$ is indeed a norm on $V$, and it will suffice to show that $\norm{\cdot},\norm{\cdot}_\infty$ are equivalent, and that $\norm{\cdot}_\infty$ is complete. For completeness, we note that a sequence of vectors is Cauchy under $\norm{\cdot}_\infty$ if and only if each of the $d$ coordinates is Cauchy, and the same goes for convergence. So let us show that $\norm{\cdot},\norm{\cdot}_\infty$ are equivalent. We have one inequality
\[\norm{v} \le (\norm{e_1}+\dots+\norm{e_n})\norm{v}_\infty\]
for all $v\in V$. Suppose, for the sake of contradiction, that we do not have a constant $c>0$ with 
\[c\norm{v}_\infty\le \norm{v}\]
for all $v\in V$.
Then, we may find a sequence of vectors $(v_n)_n$ with $\norm{v_n}_\infty = 1$ and $\norm{v_n}\to 0$. 
Restricting to a subsequence, we may assume that the $d$-th coordinate of each $v_n$ has absolute value $1$. 
Dividing by those coordinates, we may assume that the $d$-th coordinate of each $v_n$ is $1$. 
Write $v_n=e_d+u_n$ where $u_n\in\spn(e_1,\dots,e_{d-1})$. We get that $u_n$ converges to $-e_n$ under $\norm{\cdot}$. However,
by induction, the resticted norm $\norm{\cdot}$ on $\spn(e_1,\dots,e_{d-1})$ is complete, and so $\spn(e_1,\dots,e_{d-1})$ is closed in $V$ under $\norm{\cdot}$. This implies $-e_n\in \spn(e_1,\dots,e_{d-1})$, a contradiction. 
\end{proof}


\begin{fact*}
Let $\abs{\cdot}$ be a valuation on $K$. Then any point in any ball is the center of the ball.
\end{fact*}


\begin{fact*}
Let $K$ be a complete field under the absolute value $\abs{\cdot}$. Suppose $E/K$ is a finite field extension. Then there is at most one extension of the absolute value $\abs{\cdot}$ to an absolute value of $K$. 
%(TODO show there is one). the proofs in the literature (eg https://math.stanford.edu/~conrad/676Page/handouts/ostrowski.pdf) seem a bit sad, in the sense that it uses Hensel's lemma and that complete Archimedean means R/C. 
\end{fact*}



\newpage
\section{Galois Theory}
\begin{fact*}
Let $G$ be a group, and let $F$ be a field. Then, the set of group homomorphisms $G\to F^\times$ is linearly independent over $F$.
\end{fact*}


\begin{proof}
Suppose, for the sake of contradiction, that $\chi_1,\dots,\chi_d : G\to F^\times$ are distinct homomorphisms that are linearly dependent. 
We may further assume that $d$ is minimal. 
Let $a_1,\dots,a_d$ be scalars, not all zero, with \[a_1\chi_1 + \dots + a_d\chi_d = 0\]
By the minimality of $d$, we have that all the $a_i$ are non-zero. 
Now, setting $g\leftarrow hg$ we have
\[a_1\chi_1(h)\chi_1 + \dots + a_d\chi_d(h)\chi_d = 0\]
for each $h\in G$. By the minimality of $d$, the two equations must be scalar multiples of one another. Thus $\chi_1(h)=\dots=\chi_d(h)$. Since $h$ is arbitrary, $\chi_1=\dots=\chi_d$. So $d=1$ and $\chi_1=0$, which is nonsense.
\end{proof}


\begin{nomn*}
We set
$\aut(L/K) := \{\varphi: L\to L \text{ is a field isomorphism, } \varphi\restriction_{K} = \id_K\}$.
\end{nomn*}


% \begin{fact*}
% Let $L/K$ be a finite field extension of degree $d$. Then $\#\aut(L/K) \le d$.
% \end{fact*}

%\newpage
% Next lectures - 
%               - Euclidean Algorithm, gcd, lcm, continued fractions
%               - Introduction to Modular Arithmetic (chinese, fermat, euler, wilson)
%               - Introduction to d-adics
%               - Squares residues, sums of squares, quadratic reciprocity, Z[i], quadratic forms in general
%               - arithmetic functions
%               - pythagorean triples, geometric method, infinite descent
%               - misc (vieta jumping, lifting the exponent)
%               - cyclotomics
%               - finite fields, primitive elements
%               - ring of integers, factorization of ideals and class group
%               - fibonacci, chevalley warning,
%               - valuation theory, newton polygons
%               - there's a proof of Wedderburn's theorem that a finite non-commutative domain is a field using cyclotomics

%   good refs:  http://gaetan.chenevier.perso.math.cnrs.fr/GT/ireland_rosen.pdf
%               https://virtualmath1.stanford.edu/~conrad/154Page/handouts.html
%               Ben Green, Brian Conrad, Keith Conrad, De Shalit
%               https://nsnyder1.pages.iu.edu/thesismain.pdf
%               (elementary intro to langlands) https://www.ams.org/journals/bull/1984-10-02/S0273-0979-1984-15237-6/S0273-0979-1984-15237-6.pdf
%               other ideas: english translations of the books of Gauss+Dirichlet
\end{document}



%           nice proof of FTA: https://www.daniellitt.com/blog/2016/10/6/a-minimal-proof-of-the-fundamental-theorem-of-algebra
%           nice proof of weak nullstellensatz: https://blogs.mat.ucm.es/arrondo/wp-content/uploads/sites/75/2023/09/monthly169-171-arrondo.pdf